\documentclass[12pt,a4paper]{article}
\usepackage{amsmath,amsthm,amssymb,mathtools}
\usepackage{tikz}
\usetikzlibrary{arrows.meta,positioning,shapes.geometric,calc,fit,backgrounds}
\usepackage{graphicx}
\usepackage{booktabs,tabularx,multirow}
\usepackage[ruled,vlined,linesnumbered]{algorithm2e}
\usepackage[most]{tcolorbox}
\usepackage{hyperref}
\usepackage[capitalize,noabbrev]{cleveref}
\usepackage[margin=1in]{geometry}
\usepackage{fancyhdr}
\usepackage{enumitem}
\usepackage{xcolor}

%% ── tcolorbox environments ──────────────────────────────────────────────────
\newtcolorbox{keyinsight}[1][]{colback=blue!5!white,colframe=blue!75!black,fonttitle=\bfseries,title=Key Insight,#1}
\newtcolorbox{example}[1][]{colback=green!5!white,colframe=green!75!black,fonttitle=\bfseries,title=Example,#1}
\newtcolorbox{warningbox}[1][]{colback=red!5!white,colframe=red!75!black,fonttitle=\bfseries,title=Warning,#1}

%% ── amsthm environments ─────────────────────────────────────────────────────
\newtheorem{theorem}{Theorem}[section]
\newtheorem{lemma}[theorem]{Lemma}
\newtheorem{definition}{Definition}[section]
\newtheorem{remark}{Remark}[section]

%% ── header / footer ─────────────────────────────────────────────────────────
\pagestyle{fancy}
\fancyhf{}
\rhead{Session 18 -- RNN Hands-On \& LSTM Introduction}
\lhead{Deep Learning Course}
\rfoot{\thepage}

%% ── hyperref ────────────────────────────────────────────────────────────────
\hypersetup{colorlinks=true,linkcolor=blue!60!black,urlcolor=blue!60!black,
            citecolor=green!50!black}

%% ── convenience macros ──────────────────────────────────────────────────────
\newcommand{\xt}{x_t}
\newcommand{\ht}{h_t}
\newcommand{\hprev}{h_{t-1}}
\newcommand{\tanh}{\operatorname{tanh}}

%%%%%%%%%%%%%%%%%%%%%%%%%%%%%%%%%%%%%%%%%%%%%%%%%%%%%%%%%%%%%%%%%%%%%%%%%%%%%
\begin{document}
%%%%%%%%%%%%%%%%%%%%%%%%%%%%%%%%%%%%%%%%%%%%%%%%%%%%%%%%%%%%%%%%%%%%%%%%%%%%%

%% ── Title page ──────────────────────────────────────────────────────────────
\begin{titlepage}
  \centering
  \vspace*{2cm}
  {\Huge\bfseries Deep Learning Course\\[0.4em]
   Session 18 Lecture Notes\par}
  \vspace{1.2cm}
  {\LARGE\itshape RNN Hands-On Implementation \&\\[0.3em]
   Introduction to LSTM\par}
  \vspace{1.5cm}
  \rule{\linewidth}{0.6pt}\\[0.4cm]
  \begin{tabular}{ll}
    \textbf{Session Number:}  & 18 \\[4pt]
    \textbf{Topic:}           & Recurrent Neural Networks -- hands-on; \\
                              & unrolled RNN as feedforward network; \\
                              & motivation for LSTM \\[4pt]
    \textbf{Date:}            & January 17, 2026 \\[4pt]
    \textbf{Instructors:}     & Prof.\ S R M Prasanna \& Dr.\ Sunil Saumya \\[4pt]
    \textbf{Institution:}     & IIIT Dharwad \\[4pt]
    \textbf{Slide Deck:}      & \texttt{08-PyTorchDL-RNN} (slide 20/29) \\[4pt]
    \textbf{Notebook:}        & \texttt{RNNClassification.ipynb} (Google Colab) \\[4pt]
    \textbf{Dataset:}         & IMDB Movie Reviews (50\,000 samples) \\
                              & Shakespeare Dataset (500 plays) via Kaggle \\
  \end{tabular}
  \rule{\linewidth}{0.6pt}
  \vfill
  {\large These notes synthesise the lecture transcript and visual slide
   extractions for Session~18 of the deep learning course at IIIT Dharwad.
   They cover the architecture of simple RNNs, text pre-processing,
   the PyTorch implementation of a sentiment-analysis RNN, and the
   theoretical motivation for LSTM.\par}
\end{titlepage}

\tableofcontents
\newpage

%%%%%%%%%%%%%%%%%%%%%%%%%%%%%%%%%%%%%%%%%%%%%%%%%%%%%%%%%%%%%%%%%%%%%%%%%%%%%
\section{Introduction}
\label{sec:introduction}
%%%%%%%%%%%%%%%%%%%%%%%%%%%%%%%%%%%%%%%%%%%%%%%%%%%%%%%%%%%%%%%%%%%%%%%%%%%%%

Session~18 is a pivotal bridge between RNN theory and real-world
implementation.  The session opened with a brief review of L2
regularisation and a Q\&A on assignments, then devoted the majority of
its $\approx 1$h\,52\,min runtime to two interlocking themes:

\begin{enumerate}[label=\textbf{\arabic*.}]
  \item \textbf{RNN as an unrolled feedforward network.}  The instructor
        revisited the computational graph of an RNN unrolled to time step
        $t=3$, establishing the shared weight matrices $W$ and $U$, and
        deriving the hidden-state update equation.
  \item \textbf{Hands-on sentiment analysis.}  A complete PyTorch
        implementation of a simple RNN for IMDB movie-review sentiment
        classification was built from scratch inside a Google Colab
        notebook (\texttt{RNNClassification.ipynb}).  Every stage of the
        pipeline -- data loading, text pre-processing, vocabulary
        construction, dataset/dataloader classes, model definition, and
        training loop -- was coded and explained live.
\end{enumerate}

The session concluded by returning to the unrolled RNN diagram to
explain \emph{why} long sequences create a vanishing-gradient problem,
directly motivating the LSTM architecture to be introduced in the
next session.

\begin{keyinsight}
The central architectural difference between a standard ANN and an RNN
is the \emph{recurrent connection}: when processing word $x_t$, the
network also receives the hidden state $h_{t-1}$ produced when
processing $x_{t-1}$.  This allows the network to maintain a
``memory'' of previous inputs, which is essential for sequential data.
\end{keyinsight}

%%%%%%%%%%%%%%%%%%%%%%%%%%%%%%%%%%%%%%%%%%%%%%%%%%%%%%%%%%%%%%%%%%%%%%%%%%%%%
\section{Background and Prerequisites}
\label{sec:background}
%%%%%%%%%%%%%%%%%%%%%%%%%%%%%%%%%%%%%%%%%%%%%%%%%%%%%%%%%%%%%%%%%%%%%%%%%%%%%

\subsection{Sequential Data and the Limits of ANNs}
\label{subsec:ann_limits}

A standard feed-forward neural network (ANN) treats every input
independently.  Given a vocabulary of size $V$, each word $w$ is
represented as a one-hot vector $\mathbf{x} \in \{0,1\}^V$, and the
network applies:

\begin{equation}
  \mathbf{o} = f\!\left(W_{\text{in}}\,\mathbf{x} + \mathbf{b}\right),
  \label{eq:ann_forward}
\end{equation}

where $f$ is a non-linear activation function.  Two limitations arise
immediately:

\begin{enumerate}[label=(\alph*)]
  \item \textbf{Vocabulary-sized input.}  For $V = 75\,509$ (the vocabulary
        of the IMDB corpus used in this session), $W_{\text{in}}$ has
        $75\,509 \times h$ parameters for a hidden layer of width $h$.
        At $h=64$ that is already $\approx 4.8$\,M weights \emph{just}
        for the first layer.
  \item \textbf{No temporal dependency.}  Processing the words
        ``\emph{not}'' and ``\emph{good}'' simultaneously in a window
        discards the positional relationship between them.  The network
        sees both as part of a fixed-length input vector; it does not
        ``know'' which came first.
\end{enumerate}

Window-based approaches (feeding $k$ consecutive words as a single
input) partially address (b) but cannot handle variable-length
sequences without padding, and they still treat the $k$ words as a
single unordered bag.

\subsection{L2 Regularisation Review}
\label{subsec:l2_review}

The session opened with a clarification about L2 regularisation,
which had appeared in an assignment question.  The standard form adds
a penalty proportional to the squared weight magnitudes to the loss:

\begin{equation}
  \mathcal{L}_{\text{reg}} = \mathcal{L} + \frac{\lambda}{2}
  \sum_{j} w_j^2,
  \label{eq:l2_regularisation}
\end{equation}

where $\lambda$ is the regularisation strength (\texttt{weight\_decay}
in PyTorch optimisers).  The factor of $\tfrac{1}{2}$ is a
mathematical convenience: when differentiating, the $2$ in the
exponent cancels with $\tfrac{1}{2}$, yielding the clean gradient
$\nabla_{w_j}\mathcal{L}_{\text{reg}} = \nabla_{w_j}\mathcal{L} +
\lambda w_j$.

\begin{example}[title=L2 Penalty Calculation]
  Given $\lambda = 10^{-4}$ and a weight $w = 2.0$:
  \begin{align}
    \text{With factor } \tfrac{1}{2}:&\quad
      \frac{\lambda}{2} w^2 = \frac{10^{-4}}{2} \times 4 = 2 \times 10^{-4} \\
    \text{Without factor } \tfrac{1}{2}:&\quad
      \lambda w^2 = 10^{-4} \times 4 = 4 \times 10^{-4}
  \end{align}
  Both forms are mathematically valid; the choice must be consistent
  with the convention used when defining the loss function.
\end{example}

%%%%%%%%%%%%%%%%%%%%%%%%%%%%%%%%%%%%%%%%%%%%%%%%%%%%%%%%%%%%%%%%%%%%%%%%%%%%%
\section{RNN Architecture: Unrolled View}
\label{sec:rnn_architecture}
%%%%%%%%%%%%%%%%%%%%%%%%%%%%%%%%%%%%%%%%%%%%%%%%%%%%%%%%%%%%%%%%%%%%%%%%%%%%%

\subsection{Hidden-State Update Equation}
\label{subsec:hidden_state}

At every time step $t$, a simple RNN computes:

\begin{equation}
  \boxed{h_t = f\!\left(W\,h_{t-1} + U\,x_t + b\right)}
  \label{eq:rnn_update}
\end{equation}

where:
\begin{itemize}[noitemsep]
  \item $x_t \in \mathbb{R}^{d}$ is the embedding of the word at
        position $t$ ($d = 100$ in the notebook);
  \item $h_t \in \mathbb{R}^{n_h}$ is the hidden state ($n_h = 64$
        neurons);
  \item $W \in \mathbb{R}^{n_h \times n_h}$ is the
        \emph{hidden-to-hidden} weight matrix (shared across all $t$);
  \item $U \in \mathbb{R}^{n_h \times d}$ is the
        \emph{input-to-hidden} weight matrix (shared across all $t$);
  \item $b \in \mathbb{R}^{n_h}$ is the bias vector;
  \item $f$ is a point-wise non-linearity (typically $\tanh$).
\end{itemize}

The output at time step $t=3$ as shown on slide V01 is:

\begin{equation}
  O_3 = f\!\left(W \cdot O_2 + U \cdot X_3\right),
  \label{eq:o3}
\end{equation}

which is the equation explicitly displayed on the course slide.

\begin{keyinsight}[title=Parameter Sharing is the Key]
  The matrices $W$ and $U$ are \textbf{identical} at every time step.
  This is why an RNN can handle sequences of any length without
  increasing the number of parameters.  A 391-word review and a
  61-word review both use the same $W$ and $U$; only the number of
  times they are applied differs.
\end{keyinsight}

\subsection{Unrolled RNN as a Deep Feedforward Network}
\label{subsec:unrolled}

\Cref{fig:rnn_unrolled} shows the unrolled computational graph for a
sequence of length $T$.  Each column corresponds to one time step;
the rows represent layers.  The horizontal arrows represent the flow
of hidden state from $h_{t-1}$ to $h_t$, while the vertical arrows
represent the local computation within each time step.

\begin{figure}[ht]
  \centering
  \begin{tikzpicture}[
      node distance=1.8cm and 2.2cm,
      box/.style={draw,rounded corners=3pt,minimum width=1.4cm,
                  minimum height=0.8cm,fill=blue!10,font=\small},
      inp/.style={draw,circle,minimum size=0.7cm,fill=orange!20,font=\small},
      out/.style={draw,circle,minimum size=0.7cm,fill=green!20,font=\small},
      arr/.style={-{Stealth[length=5pt]},thick},
      rarr/.style={-{Stealth[length=5pt]},thick,blue!70!black}
    ]

    %% Hidden state nodes
    \node[box] (h0) {$h_0$};
    \node[box,right=of h0] (h1) {$h_1$};
    \node[box,right=of h1] (h2) {$h_2$};
    \node[box,right=of h2] (h3) {$h_3$};

    %% Recurrent arrows
    \draw[rarr] (h0) -- node[above,font=\tiny]{$W$} (h1);
    \draw[rarr] (h1) -- node[above,font=\tiny]{$W$} (h2);
    \draw[rarr] (h2) -- node[above,font=\tiny]{$W$} (h3);

    %% Input nodes (below)
    \node[inp,below=1.2cm of h0] (x0) {$x_0$};
    \node[inp,below=1.2cm of h1] (x1) {$x_1$};
    \node[inp,below=1.2cm of h2] (x2) {$x_2$};
    \node[inp,below=1.2cm of h3] (x3) {$x_3$};

    %% Input arrows
    \draw[arr] (x0) -- node[right,font=\tiny]{$U$} (h0);
    \draw[arr] (x1) -- node[right,font=\tiny]{$U$} (h1);
    \draw[arr] (x2) -- node[right,font=\tiny]{$U$} (h2);
    \draw[arr] (x3) -- node[right,font=\tiny]{$U$} (h3);

    %% Output nodes (above)
    \node[out,above=1.2cm of h0] (o0) {$O_0$};
    \node[out,above=1.2cm of h1] (o1) {$O_1$};
    \node[out,above=1.2cm of h2] (o2) {$O_2$};
    \node[out,above=1.2cm of h3] (o3) {$O_3$};

    %% Output arrows
    \draw[arr] (h0) -- (o0);
    \draw[arr] (h1) -- (o1);
    \draw[arr] (h2) -- (o2);
    \draw[arr] (h3) -- (o3);

    %% Labels
    \node[font=\footnotesize,below=0.2cm of x0] {$t=0$};
    \node[font=\footnotesize,below=0.2cm of x1] {$t=1$};
    \node[font=\footnotesize,below=0.2cm of x2] {$t=2$};
    \node[font=\footnotesize,below=0.2cm of x3] {$t=3$};

    %% Initial state h_{-1}
    \node[left=1.2cm of h0,font=\small] (hinit) {$h_{-1}=\mathbf{0}$};
    \draw[rarr] (hinit) -- node[above,font=\tiny]{$W$} (h0);

  \end{tikzpicture}
  \caption{RNN unrolled over four time steps $t = 0, 1, 2, 3$.  The
           weight matrices $W$ (hidden-to-hidden) and $U$
           (input-to-hidden) are \emph{shared} at every column.
           The highlighted output $O_3$ at $t=3$ satisfies
           $O_3 = f(W \cdot O_2 + U \cdot X_3)$, as shown on slide V01.}
  \label{fig:rnn_unrolled}
\end{figure}

\begin{remark}
  The unrolled view makes clear that an RNN with sequence length $T$
  is equivalent to a feedforward network of depth $T$.  For a 391-word
  IMDB review, the unrolled network has 391 ``layers'', all sharing
  the same parameters $W$, $U$.  This depth is precisely why
  gradients can vanish during backpropagation through time (BPTT).
\end{remark}

\subsection{ANN vs.\ RNN: Architectural Comparison}
\label{subsec:ann_vs_rnn}

\Cref{tab:ann_vs_rnn} summarises the key differences between a
standard ANN and an RNN for sequential NLP tasks, as discussed by
the instructor during the session.

\begin{table}[ht]
  \centering
  \caption{ANN versus RNN for sequential text processing.}
  \label{tab:ann_vs_rnn}
  \begin{tabularx}{\linewidth}{lXX}
    \toprule
    \textbf{Property} & \textbf{ANN} & \textbf{RNN} \\
    \midrule
    Input at time $t$ & Fixed-size window of $k$ words
                      & Single word $x_t$; also receives $h_{t-1}$ \\
    \addlinespace
    Sequence order & Implicit (position in window)
                   & Explicit (strictly sequential) \\
    \addlinespace
    Variable-length input & Requires padding to fixed $k$
                          & Handled naturally (no padding needed) \\
    \addlinespace
    Parameter count & $V \times h$ for first layer ($V = 75\,509$)
                    & $d \times n_h + n_h \times n_h$ (after embedding) \\
    \addlinespace
    Long-range context & Limited to window size $k$
                       & Theoretically unlimited (in practice limited by vanishing gradient) \\
    \addlinespace
    Parallelism & All words in window processed simultaneously
                & Sequential -- $t+1$ requires $h_t$ \\
    \addlinespace
    Typical use & Classification with fixed context
                & Sequence classification, language modelling \\
    \bottomrule
  \end{tabularx}
\end{table}

%%%%%%%%%%%%%%%%%%%%%%%%%%%%%%%%%%%%%%%%%%%%%%%%%%%%%%%%%%%%%%%%%%%%%%%%%%%%%
\section{Text Pre-Processing Pipeline}
\label{sec:preprocessing}
%%%%%%%%%%%%%%%%%%%%%%%%%%%%%%%%%%%%%%%%%%%%%%%%%%%%%%%%%%%%%%%%%%%%%%%%%%%%%

The IMDB dataset contains 50\,000 labelled movie reviews (25\,000
positive, 25\,000 negative).  The session used a randomly sampled
subset of 5\,000 reviews to keep training times manageable on Colab.
\Cref{fig:preprocessing_pipeline} illustrates the five-stage pipeline.

\begin{figure}[ht]
  \centering
  \begin{tikzpicture}[
      stage/.style={draw,rectangle,rounded corners=4pt,
                    minimum width=2.4cm,minimum height=0.9cm,
                    fill=teal!15,font=\small,align=center,thick},
      arr/.style={-{Stealth[length=6pt]},thick,teal!60!black}
    ]
    \node[stage] (s1) {Raw\\Text};
    \node[stage,right=1.3cm of s1] (s2) {Lower-\\case};
    \node[stage,right=1.3cm of s2] (s3) {Remove\\Punct.};
    \node[stage,right=1.3cm of s3] (s4) {Tokenise};
    \node[stage,right=1.3cm of s4] (s5) {Build\\Vocab};
    \node[stage,below=1.2cm of s5] (s6) {Text\\$\to$ Indices};
    \node[stage,below=1.2cm of s3] (s7) {Tensor\\Dataset};
    \node[stage,below=1.2cm of s1] (s8) {Data\\Loader};

    \draw[arr] (s1) -- (s2);
    \draw[arr] (s2) -- (s3);
    \draw[arr] (s3) -- (s4);
    \draw[arr] (s4) -- (s5);
    \draw[arr] (s5) -- (s6);
    \draw[arr] (s6) -- (s7);
    \draw[arr] (s7) -- (s8);

    %% labels below arrows
    \node[font=\tiny,below=0.1cm of s2] {step 1};
    \node[font=\tiny,below=0.1cm of s3] {step 2};
    \node[font=\tiny,below=0.1cm of s4] {step 3};
    \node[font=\tiny,below=0.1cm of s5] {step 4};
  \end{tikzpicture}
  \caption{Five-stage text pre-processing pipeline implemented in
           \texttt{RNNClassification.ipynb}.}
  \label{fig:preprocessing_pipeline}
\end{figure}

\subsection{Tokenisation and Cleaning}
\label{subsec:tokenisation}

The \texttt{tokenize(text)} function applied the following
transformations in sequence:

\begin{enumerate}[label=\arabic*.]
  \item Convert to lower-case.
  \item Remove all punctuation characters using the \texttt{string.punctuation} set.
  \item Split on whitespace to produce a list of tokens.
\end{enumerate}

These three steps are illustrated below in Python pseudocode:

\begin{verbatim}
import string

def tokenize(text):
    text = text.lower()
    text = ''.join(c for c in text if c not in string.punctuation)
    tokens = text.split()
    return tokens
\end{verbatim}

\begin{example}[title=Tokenisation in Action]
  Input: \texttt{``One of the most beautifully crafted films you'll ever see.''}\\[4pt]
  After lower-case: \texttt{``one of the most beautifully crafted films youll ever see''}\\[4pt]
  After removing punctuation: same (apostrophe removed; ``you'll'' $\to$ ``youll'')\\[4pt]
  After splitting: \texttt{[`one', `of', `the', `most', `beautifully', `crafted', `films', `youll', `ever', `see']}
\end{example}

Note that the apostrophe in ``you'll'' is stripped, making the
token \texttt{youll} (index~190 in the vocabulary).  A student asked
why this was not classified as the \texttt{<UNK>} token -- the answer
is that tokenisation removes the apostrophe \emph{before} vocabulary
lookup, so \texttt{youll} appears as a known token in the vocabulary.

\subsection{Vocabulary Construction}
\label{subsec:vocabulary}

The vocabulary is a Python dictionary mapping each unique token to an
integer index:

\begin{equation}
  \text{vocab} = \bigl\{\,\text{token} \mapsto \text{index}\,\bigr\}
  \label{eq:vocab}
\end{equation}

Two special tokens are pre-inserted before any text is processed:

\begin{center}
  \begin{tabular}{cll}
    \toprule
    \textbf{Index} & \textbf{Token} & \textbf{Purpose} \\
    \midrule
    0 & \texttt{<PAD>} & Padding (needed when batches contain variable-length sequences) \\
    1 & \texttt{<UNK>} & Unknown token (used at inference for out-of-vocabulary words) \\
    \bottomrule
  \end{tabular}
\end{center}

Every new unique token encountered while iterating through the 5\,000
training reviews is assigned the next available index
($|\text{vocab}|$ at the time of insertion):

\begin{verbatim}
def build_vocab(texts):
    vocab = {'<PAD>': 0, '<UNK>': 1}
    for text in texts:
        for token in tokenize(text):
            if token not in vocab:
                vocab[token] = len(vocab)
    return vocab
\end{verbatim}

After processing all 5\,000 rows the vocabulary contained
$|\text{vocab}| = 75\,509$ unique tokens.  The instructor noted that
processing the full 50\,000-sample dataset would likely exceed
100\,000 unique tokens.

\begin{keyinsight}[title=Why Both \texttt{<PAD>} and \texttt{<UNK>}?]
  \textbf{\texttt{<UNK>}} (index 1) handles out-of-vocabulary words at
  test time.  If a review contains a word not seen during training, it
  is mapped to index~1.  \textbf{\texttt{<PAD>}} (index 0) is included
  as a standard convention even though this notebook processes each
  review individually (batch size = 1), so no padding is actually
  applied.  Padding becomes necessary when batching multiple
  variable-length sequences together.
\end{keyinsight}

\subsection{Text-to-Indices Conversion}
\label{subsec:text_to_indices}

The function \texttt{text\_to\_indices(text, vocab)} converts a raw
review string into a tensor of vocabulary indices:

\begin{verbatim}
def text_to_indices(text, vocab):
    tokens = tokenize(text)
    return [vocab[t] if t in vocab else vocab['<UNK>'] for t in tokens]
\end{verbatim}

The sequence of indices preserves the word order of the original
review.  For example, the sentence ``one of the \ldots'' maps to the
integer sequence $[162, 12, 16, \ldots]$.  This is the input
$[x_0, x_1, x_2, \ldots]$ that will be fed token-by-token to the RNN.

\subsection{Padding and Batch Processing Trade-offs}
\label{subsec:padding}

A student raised the question of padding during the live session.
\Cref{tab:padding_tradeoffs} summarises the trade-offs discussed:

\begin{table}[ht]
  \centering
  \caption{Trade-offs between padded batching and individual sequence processing.}
  \label{tab:padding_tradeoffs}
  \begin{tabularx}{\linewidth}{lXX}
    \toprule
    \textbf{Property} & \textbf{Batch with Padding} & \textbf{Batch Size = 1} \\
    \midrule
    GPU utilisation & Higher (parallelism) & Lower \\
    \addlinespace
    Information loss & Yes -- zeros dilute gradient signal & None \\
    \addlinespace
    Memory & Fixed (max-length sequences) & Variable \\
    \addlinespace
    Implementation & Requires \texttt{<PAD>} masking & Simple \\
    \addlinespace
    Recommended for RNN & Only when resource-constrained & Preferred \\
    \bottomrule
  \end{tabularx}
\end{table}

\begin{warningbox}
  Padding with zeros is never ideal for RNNs.  The padded time steps
  produce meaningless hidden-state updates, dilute the gradient signal
  during backpropagation, and can degrade model performance.  The
  best practice is to process each sequence individually
  (\texttt{batch\_size = 1}) or to use \texttt{pack\_padded\_sequence}
  in PyTorch, which skips padded positions entirely.
\end{warningbox}

%%%%%%%%%%%%%%%%%%%%%%%%%%%%%%%%%%%%%%%%%%%%%%%%%%%%%%%%%%%%%%%%%%%%%%%%%%%%%
\section{PyTorch Dataset and DataLoader}
\label{sec:dataset_dataloader}
%%%%%%%%%%%%%%%%%%%%%%%%%%%%%%%%%%%%%%%%%%%%%%%%%%%%%%%%%%%%%%%%%%%%%%%%%%%%%

\subsection{Custom Dataset Class}
\label{subsec:custom_dataset}

Following the standard PyTorch pattern introduced in earlier sessions,
a custom \texttt{IMDBDataset} class was defined by inheriting from
\texttt{torch.utils.data.Dataset} and implementing three methods:

\begin{itemize}[noitemsep]
  \item \texttt{\_\_init\_\_(self, df, vocab)}: Stores the DataFrame
        (with a reset index) and the vocabulary dictionary.
  \item \texttt{\_\_len\_\_(self)}: Returns the number of reviews.
  \item \texttt{\_\_getitem\_\_(self, idx)}: Converts the $idx$-th
        review to a tensor of indices and returns the
        (review\_tensor, label\_tensor) pair.
\end{itemize}

The \texttt{\_\_getitem\_\_} method calls \texttt{text\_to\_indices}
on-the-fly, so no pre-encoded copy of the entire dataset needs to be
held in memory.  The label is returned as a floating-point tensor
($0.0$ for negative, $1.0$ for positive) to match the binary
cross-entropy loss expectation.

\begin{verbatim}
class IMDBDataset(Dataset):
    def __init__(self, df, vocab):
        self.df = df.reset_index(drop=True)
        self.vocab = vocab

    def __len__(self):
        return len(self.df)

    def __getitem__(self, idx):
        review = self.df.loc[idx, 'review']
        label  = self.df.loc[idx, 'sentiment']
        x = torch.tensor(text_to_indices(review, self.vocab))
        y = torch.tensor(float(label))
        return x, y
\end{verbatim}

The reason for \texttt{reset\_index(drop=True)} is that the 5\,000
rows were drawn by \emph{random sampling} from the original 50\,000,
so the index values are non-contiguous (e.g.\ 4, 312, 7003, \ldots).
Resetting creates a clean 0-to-4999 integer index that \texttt{iloc}
can use directly.

\subsection{DataLoader Configuration}
\label{subsec:dataloader}

The DataLoader was configured with \texttt{batch\_size = 1} for the
reasons discussed in \cref{subsec:padding}:

\begin{verbatim}
train_loader = DataLoader(train_dataset,
                          batch_size=1,
                          shuffle=True)
\end{verbatim}

Iterating the loader yields batches of shape
$1 \times T_i$ for the review tensor (where $T_i$ is the number of
tokens in the $i$-th review) and shape $1$ for the label scalar.  The
instructor demonstrated this live, showing:

\begin{itemize}[noitemsep]
  \item Batch 0: review shape $1 \times 61$ (61 tokens) -- label: positive.
  \item Batch 1: review shape $1 \times 391$ (391 tokens) -- label: negative.
  \item Batch 2: review shape $1 \times 173$ (173 tokens).
\end{itemize}

This confirms that RNNs handle variable-length inputs naturally -- the
DataLoader does not need to pad or truncate any sequence.

%%%%%%%%%%%%%%%%%%%%%%%%%%%%%%%%%%%%%%%%%%%%%%%%%%%%%%%%%%%%%%%%%%%%%%%%%%%%%
\section{Simple RNN Model Implementation}
\label{sec:model}
%%%%%%%%%%%%%%%%%%%%%%%%%%%%%%%%%%%%%%%%%%%%%%%%%%%%%%%%%%%%%%%%%%%%%%%%%%%%%

\subsection{Model Architecture Overview}
\label{subsec:model_overview}

The \texttt{SimpleRNN} model consists of exactly three learnable
components stacked in sequence:

\begin{enumerate}[label=\arabic*.]
  \item \textbf{Embedding layer} (\texttt{nn.Embedding}): maps each
        integer token index to a dense $d$-dimensional vector.
  \item \textbf{RNN layer} (\texttt{nn.RNN}): processes the embedded
        sequence token-by-token, maintaining a hidden state.
  \item \textbf{Linear (fully-connected) layer} (\texttt{nn.Linear}):
        maps the final hidden state to a scalar sentiment score.
\end{enumerate}

\Cref{fig:model_architecture} illustrates the full forward pass.

\begin{figure}[ht]
  \centering
  \begin{tikzpicture}[
      layer/.style={draw,rectangle,rounded corners=5pt,
                    minimum width=3.5cm,minimum height=0.9cm,
                    fill=violet!12,font=\small,align=center,thick},
      data/.style={draw,rectangle,rounded corners=3pt,
                   minimum width=3.5cm,minimum height=0.7cm,
                   fill=gray!12,font=\small,align=center},
      arr/.style={-{Stealth[length=6pt]},thick}
    ]

    \node[data]  (inp)   {Token indices $[i_0,i_1,\ldots,i_{T-1}]$};
    \node[layer,below=0.8cm of inp]  (emb)
        {Embedding\\$V \times 100 \to T \times 100$};
    \node[layer,below=0.8cm of emb]  (rnn)
        {RNN\\input=100, hidden=64\\(batch\_first=True)};
    \node[data,below=0.6cm of rnn]   (hn)
        {$h_T \in \mathbb{R}^{64}$ (final hidden state)};
    \node[layer,below=0.6cm of hn]   (fc)
        {Linear $64 \to 1$};
    \node[data,below=0.6cm of fc]    (out)
        {Scalar logit $\hat{y} \in \mathbb{R}$};
    \node[layer,below=0.6cm of out,fill=orange!15]  (loss)
        {BCEWithLogitsLoss\\$\mathcal{L}(\hat{y}, y)$};

    \draw[arr] (inp)  -- (emb);
    \draw[arr] (emb)  -- node[right,font=\tiny]{embedded sequence} (rnn);
    \draw[arr] (rnn)  -- node[right,font=\tiny]{output, $h_n$} (hn);
    \draw[arr] (hn)   -- (fc);
    \draw[arr] (fc)   -- (out);
    \draw[arr] (out)  -- (loss);

  \end{tikzpicture}
  \caption{Forward pass of the \texttt{SimpleRNN} model used in the
           session.  $V = 75\,509$ (vocabulary size), $T$ = sequence
           length (variable), $d = 100$ (embedding dimension),
           $n_h = 64$ (RNN hidden size).}
  \label{fig:model_architecture}
\end{figure}

\subsection{The Embedding Layer}
\label{subsec:embedding}

The embedding layer performs two roles simultaneously:

\begin{enumerate}[label=(\alph*)]
  \item \textbf{One-hot encoding:} internally, the layer treats each
        integer index as selecting the corresponding row from a learnable
        weight matrix $E \in \mathbb{R}^{V \times d}$.  This is
        mathematically identical to multiplying a one-hot vector by $E$.
  \item \textbf{Dimensionality reduction:} the 75\,509-dimensional
        one-hot vector is compressed to a dense $d=100$-dimensional
        vector.
\end{enumerate}

\begin{equation}
  \mathbf{e}_t = E[i_t] \in \mathbb{R}^{d}, \quad d = 100
  \label{eq:embedding}
\end{equation}

The two key advantages of using embeddings rather than raw one-hot
vectors:

\begin{enumerate}[noitemsep]
  \item \textbf{Parameter reduction:} the input-to-hidden weight
        matrix changes from $V \times n_h = 75\,509 \times 64
        \approx 4.8$\,M parameters to $d \times n_h = 100 \times 64
        = 6\,400$ parameters -- a reduction by a factor of $\approx 750$.
  \item \textbf{Denser gradients:} a one-hot vector with 75\,508 zeros
        produces almost no gradient signal.  A dense embedding vector
        with non-zero values everywhere propagates gradients more
        effectively.
\end{enumerate}

\begin{keyinsight}[title=Embedding Layer as Fully-Connected Layer]
  Conceptually, the embedding layer is simply a fully-connected layer
  (without bias or activation) placed \emph{before} the main
  architecture starts.  Its weights $E$ are initialised randomly and
  learned by gradient descent along with all other parameters.  If
  pre-trained word embeddings (Word2Vec, GloVe, FastText) are loaded
  into $E$, the model can leverage semantic information from large
  external corpora.
\end{keyinsight}

\subsection{The RNN Layer}
\label{subsec:rnn_layer}

The \texttt{nn.RNN} module implements \cref{eq:rnn_update} for every
time step in a vectorised manner.  Key constructor arguments used:

\begin{center}
  \begin{tabular}{lll}
    \toprule
    \textbf{Argument} & \textbf{Value} & \textbf{Meaning} \\
    \midrule
    \texttt{input\_size}  & 100  & Dimension of each embedded token \\
    \texttt{hidden\_size} & 64   & Width of the hidden state $h_t$ \\
    \texttt{batch\_first} & \texttt{True} & Input tensor shape: $(B, T, d)$ not $(T, B, d)$ \\
    \bottomrule
  \end{tabular}
\end{center}

When called, \texttt{nn.RNN} returns two tensors:
\begin{itemize}[noitemsep]
  \item \texttt{output}: shape $(B, T, n_h)$ -- the hidden state at
        \emph{every} time step.
  \item \texttt{h\_n}: shape $(1, B, n_h)$ -- the hidden state at the
        \emph{last} time step only.
\end{itemize}

For sentiment classification, only \texttt{h\_n} is needed because
the final hidden state encodes a summary of the entire sequence.
The instructor verified this by showing that the last row of
\texttt{output} is numerically identical to \texttt{h\_n}.

\subsection{Parameter Count}
\label{subsec:param_count}

The total number of trainable parameters in the model is:

\begin{align}
  \text{Embedding:} &\quad V \times d = 75\,509 \times 100 = 7\,550\,900 \label{eq:emb_params} \\
  \text{RNN} \;(U): &\quad d \times n_h = 100 \times 64 = 6\,400 \label{eq:rnn_u} \\
  \text{RNN} \;(W): &\quad n_h \times n_h = 64 \times 64 = 4\,096 \label{eq:rnn_w} \\
  \text{RNN biases:}&\quad 2 \times n_h = 128 \label{eq:rnn_bias} \\
  \text{Linear:}    &\quad n_h \times 1 + 1 = 65 \label{eq:fc_params} \\
  \hline \nonumber \\
  \text{Total:}     &\quad \approx 7\,561\,589 \label{eq:total_params}
\end{align}

The embedding layer dominates the parameter count.  This is typical
for NLP models with large vocabularies; pre-training embeddings or
using sub-word tokenisation (BPE) are common mitigations.

\subsection{Forward Pass Code}
\label{subsec:forward_pass}

\begin{verbatim}
class SimpleRNN(nn.Module):
    def __init__(self, vocab_size, embed_dim, hidden_size, output_size):
        super(SimpleRNN, self).__init__()
        self.embedding = nn.Embedding(vocab_size, embed_dim)
        self.rnn = nn.RNN(embed_dim, hidden_size, batch_first=True)
        self.fc  = nn.Linear(hidden_size, output_size)

    def forward(self, x):
        embedded = self.embedding(x)          # (B, T, 100)
        _, h_n   = self.rnn(embedded)         # h_n: (1, B, 64)
        h_n      = h_n.squeeze(0)             # (B, 64)
        out      = self.fc(h_n)               # (B, 1)
        return out
\end{verbatim}

The underscore \texttt{\_} discards the per-timestep output tensor;
only the final hidden state \texttt{h\_n} is passed to the
fully-connected layer.

%%%%%%%%%%%%%%%%%%%%%%%%%%%%%%%%%%%%%%%%%%%%%%%%%%%%%%%%%%%%%%%%%%%%%%%%%%%%%
\section{Training Procedure}
\label{sec:training}
%%%%%%%%%%%%%%%%%%%%%%%%%%%%%%%%%%%%%%%%%%%%%%%%%%%%%%%%%%%%%%%%%%%%%%%%%%%%%

\subsection{Loss Function and Optimiser}
\label{subsec:loss_optimiser}

Binary cross-entropy with logits was used as the loss function:

\begin{equation}
  \mathcal{L} = -\frac{1}{N}\sum_{i=1}^{N}
    \Bigl[ y_i \log \sigma(\hat{y}_i)
         + (1 - y_i) \log (1 - \sigma(\hat{y}_i)) \Bigr],
  \label{eq:bce_loss}
\end{equation}

where $\sigma$ is the sigmoid function.  \texttt{BCEWithLogitsLoss}
applies the sigmoid internally, which is numerically more stable than
applying it explicitly and then computing the log.

A student asked why a ``linear'' output layer is used instead of
logistic regression.  The answer: the linear layer produces a logit
$\hat{y}$; the \emph{sigmoid} non-linearity is applied inside
\texttt{BCEWithLogitsLoss}, so the full model does implement logistic
regression -- the non-linearity just appears in the loss function rather
than in \texttt{forward()}.

\subsection{Training Loop}
\label{subsec:training_loop}

\Cref{alg:training_loop} presents the training algorithm used in the
session.

\begin{algorithm}[ht]
  \caption{RNN Sentiment Classifier Training Loop}
  \label{alg:training_loop}
  \SetAlgoLined
  \KwIn{Training DataLoader \texttt{train\_loader}, number of epochs
        \texttt{num\_epochs}, model, loss function \texttt{criterion},
        optimiser \texttt{optimizer}}
  \KwOut{Trained model parameters}
  \BlankLine
  model.train()\;
  \For{epoch $\leftarrow 1$ \KwTo \texttt{num\_epochs}}{
    epoch\_loss $\leftarrow 0$\;
    \For{(review, label) \textbf{in} train\_loader}{
      review, label $\leftarrow$ review.to(device), label.to(device)\;
      optimizer.zero\_grad()\;
      output $\leftarrow$ model(review)\tcp*{forward pass}
      output $\leftarrow$ output.squeeze(1)\tcp*{match label shape}
      loss $\leftarrow$ criterion(output, label)\tcp*{compute BCE loss}
      loss.backward()\tcp*{backprop through time (BPTT)}
      optimizer.step()\tcp*{SGD / Adam weight update}
      epoch\_loss $\leftarrow$ epoch\_loss $+$ loss.item()\;
    }
    \textbf{print} epoch, epoch\_loss / len(train\_loader)\;
  }
\end{algorithm}

\subsection{Key Training Observations}
\label{subsec:training_observations}

\begin{itemize}
  \item The model was trained for 10 epochs on the 5\,000-sample subset.
  \item The instructor showed the loss decreasing epoch by epoch,
        confirming that the pipeline was correctly connected.
  \item Because \texttt{batch\_size = 1}, each gradient step is computed
        from a single review -- this is stochastic gradient descent
        (SGD) in its purest form, with very noisy gradient estimates.
  \item Moving to GPU (\texttt{device = 'cuda'} if available) was
        highlighted as essential for larger datasets or more epochs.
\end{itemize}

%%%%%%%%%%%%%%%%%%%%%%%%%%%%%%%%%%%%%%%%%%%%%%%%%%%%%%%%%%%%%%%%%%%%%%%%%%%%%
\section{Hands-On: Shakespeare Next-Character Prediction}
\label{sec:shakespeare}
%%%%%%%%%%%%%%%%%%%%%%%%%%%%%%%%%%%%%%%%%%%%%%%%%%%%%%%%%%%%%%%%%%%%%%%%%%%%%

The visual extractions (V02--V04) also captured code from a
secondary notebook focused on \emph{next-character prediction} using
the Kaggle Shakespeare dataset (500 plays).  This task trains an RNN
to predict the next character $c_{t+1}$ given a sequence of
characters $[c_1, c_2, \ldots, c_t]$.

\subsection{Dataset Download}
\label{subsec:shakespeare_data}

\begin{verbatim}
import kagglehub
path = kagglehub.dataset_download(
    "kylehobbsdev/shakespeare-dataset-of-500-shakespeare-plays"
)
print("Path to dataset files:", path)
# Path to dataset files:
#   /root/.cache/kaggle/datasets/kylehobbsdev/
#   shakespeare-dataset-of-500-shakespeare-plays/versions/1
\end{verbatim}

\subsection{Sliding-Window Sequence Construction}
\label{subsec:sliding_window}

Given the encoded Shakespeare text $[c_1, c_2, \ldots, c_N]$ (where
each $c_i$ is a character index), training pairs are constructed as:

\begin{equation}
  (x^{(k)},\, y^{(k)}) = \bigl([c_k, c_{k+1}, \ldots, c_{k+L-1}],\;
                                 c_{k+L}\bigr), \quad
  k = 1, 2, \ldots, N - L
  \label{eq:sliding_window}
\end{equation}

where $L$ is the context length.  The session notebook constructed
942 such training sequences from the sampled Shakespeare corpus.

\begin{example}[title=Sliding Window Sequences]
  For an encoded text starting with character indices $[3, 7, 5, 2, \ldots]$
  with context length $L = 2$:
  \begin{align*}
    (x^{(1)}, y^{(1)}) &= ([3, 7],\; 5) \\
    (x^{(2)}, y^{(2)}) &= ([3, 7, 5],\; 2) \\
    (x^{(3)}, y^{(3)}) &= ([3, 7, 5, 2],\; \ldots)
  \end{align*}
\end{example}

%%%%%%%%%%%%%%%%%%%%%%%%%%%%%%%%%%%%%%%%%%%%%%%%%%%%%%%%%%%%%%%%%%%%%%%%%%%%%
\section{Vanishing Gradient and LSTM Motivation}
\label{sec:lstm_motivation}
%%%%%%%%%%%%%%%%%%%%%%%%%%%%%%%%%%%%%%%%%%%%%%%%%%%%%%%%%%%%%%%%%%%%%%%%%%%%%

\subsection{Backpropagation Through Time (BPTT)}
\label{subsec:bptt}

Training an RNN requires computing gradients of the loss
$\mathcal{L}$ with respect to the shared weights $W$ and $U$.
Because the hidden state at time $t$ depends on all previous hidden
states, the chain rule produces products of the Jacobian matrices of
$f$:

\begin{equation}
  \frac{\partial \mathcal{L}}{\partial W}
    = \sum_{t=1}^{T} \frac{\partial \mathcal{L}}{\partial h_T}
      \underbrace{\prod_{k=t}^{T-1}
        \frac{\partial h_{k+1}}{\partial h_k}}_{\text{product of }T-t\text{ Jacobians}}
      \frac{\partial h_t}{\partial W}.
  \label{eq:bptt}
\end{equation}

If the spectral norm of each Jacobian $\tfrac{\partial h_{k+1}}{\partial h_k}$
is less than 1 (which commonly occurs with $\tanh$ activations when
$\|W\| < 1$), then for large $T$:

\begin{equation}
  \prod_{k=t}^{T-1} \left\| \frac{\partial h_{k+1}}{\partial h_k} \right\|
  \;\longrightarrow\; 0 \quad \text{as } T - t \to \infty.
  \label{eq:vanishing}
\end{equation}

This is the \textbf{vanishing gradient problem}: gradients for early
time steps become negligibly small, making it impossible to learn
long-range dependencies.

\begin{keyinsight}[title=Unrolled View Explains Vanishing Gradients]
  The unrolled RNN (slide V05) has \emph{depth} equal to sequence
  length $T$.  A 391-word review produces an unrolled network 391
  layers deep.  Just as very deep feedforward networks suffer from
  vanishing gradients in their early layers, very long RNN sequences
  lose gradient signal for early time steps.  LSTM (Long Short-Term
  Memory) solves this by introducing a \emph{cell state} and
  \emph{gating mechanisms} that allow gradients to flow through long
  sequences without vanishing.
\end{keyinsight}

\subsection{Properties of Long Sequences in RNNs}
\label{subsec:long_seq}

Two further issues arise with long sequences:

\begin{enumerate}[label=\arabic*.]
  \item \textbf{Slow processing.}  Because computation is strictly
        sequential (time step $t+1$ waits for $t$), a 500-word review
        requires 500 sequential matrix-vector multiplications.
        Compare this to a Transformer, which processes all positions
        in parallel.
  \item \textbf{Context degradation.}  Even without vanishing
        gradients, the hidden state $h_t$ is a fixed-size vector
        ($n_h = 64$) that must encode the entire history
        $[x_0, \ldots, x_{t-1}]$.  For very long sequences, the
        ``memory'' of the earliest tokens is overwritten by later ones.
\end{enumerate}

The instructor addressed a student's question about building a
question-answering system from a 100-page PDF using RNNs:

\begin{itemize}
  \item Possible in principle: convert the PDF to a structured dataset
        (e.g.\ one row per paragraph), then train a sequence-to-sequence
        RNN.
  \item Not ideal: RNNs struggle with the very long context of an entire
        book, and modern systems (RAG + large language models) perform
        significantly better on such tasks.
\end{itemize}

%%%%%%%%%%%%%%%%%%%%%%%%%%%%%%%%%%%%%%%%%%%%%%%%%%%%%%%%%%%%%%%%%%%%%%%%%%%%%
\section{Comparison of RNN Variants}
\label{sec:rnn_variants}
%%%%%%%%%%%%%%%%%%%%%%%%%%%%%%%%%%%%%%%%%%%%%%%%%%%%%%%%%%%%%%%%%%%%%%%%%%%%%

\Cref{tab:rnn_variants} compares the vanilla RNN introduced in this
session with the more advanced architectures (LSTM and GRU) mentioned
as upcoming topics.

\begin{table}[ht]
  \centering
  \caption{Comparison of RNN, LSTM, and GRU architectures.}
  \label{tab:rnn_variants}
  \begin{tabularx}{\linewidth}{lXXX}
    \toprule
    \textbf{Property} & \textbf{Vanilla RNN} & \textbf{LSTM} & \textbf{GRU} \\
    \midrule
    Hidden state & $h_t$ only & $h_t$ + cell state $c_t$ & $h_t$ only \\
    \addlinespace
    Gates & None & Input, forget, output (3) & Reset, update (2) \\
    \addlinespace
    Vanishing gradient & Severe for long $T$ & Mitigated by cell state & Partially mitigated \\
    \addlinespace
    Parameters (per layer) & $d n_h + n_h^2 + 2n_h$ & $4(d n_h + n_h^2 + 2n_h)$ & $3(d n_h + n_h^2 + 2n_h)$ \\
    \addlinespace
    Training speed & Fastest & Slowest & Intermediate \\
    \addlinespace
    Long-range dependency & Poor & Good & Good \\
    \addlinespace
    Best for & Short sequences & Long sequences & Long sequences \\
    \bottomrule
  \end{tabularx}
\end{table}

\begin{remark}
  For the IMDB sentiment task, both LSTM and GRU typically achieve
  better accuracy than vanilla RNN because movie reviews can be
  hundreds of words long and sentiment often depends on long-range
  context (e.g.\ a negation at the start of the review).
\end{remark}

%%%%%%%%%%%%%%%%%%%%%%%%%%%%%%%%%%%%%%%%%%%%%%%%%%%%%%%%%%%%%%%%%%%%%%%%%%%%%
\section{Summary and Conclusion}
\label{sec:summary}
%%%%%%%%%%%%%%%%%%%%%%%%%%%%%%%%%%%%%%%%%%%%%%%%%%%%%%%%%%%%%%%%%%%%%%%%%%%%%

\subsection{What Was Covered}
\label{subsec:covered}

Session~18 covered the following material:

\begin{enumerate}[label=\textbf{\arabic*.}]
  \item \textbf{L2 Regularisation Review.}  Both forms
        ($\lambda \sum w^2$ and $\tfrac{\lambda}{2} \sum w^2$) are
        mathematically valid; the $\tfrac{1}{2}$ factor is a
        convenience that simplifies the gradient.

  \item \textbf{RNN Architectural Motivation.}  The key change from
        ANN to RNN is the recurrent connection that feeds the previous
        hidden state $h_{t-1}$ as an additional input at each time
        step $t$.  This gives the network an explicit ``memory'' of the
        sequence so far.

  \item \textbf{Text Pre-Processing Pipeline.}  A five-stage pipeline
        was implemented: lower-casing, punctuation removal,
        tokenisation, vocabulary construction (with \texttt{<PAD>} and
        \texttt{<UNK>}), and text-to-index mapping.  The IMDB
        5\,000-sample vocabulary contained 75\,509 unique tokens.

  \item \textbf{PyTorch Dataset and DataLoader.}  A custom
        \texttt{IMDBDataset} class was created following the standard
        PyTorch pattern, with \texttt{batch\_size = 1} to avoid padding.

  \item \textbf{SimpleRNN Model.}  The model comprises three layers:
        \texttt{nn.Embedding}($V$, 100),
        \texttt{nn.RNN}(100, 64, \texttt{batch\_first=True}), and
        \texttt{nn.Linear}(64, 1).  The embedding layer converts sparse
        one-hot vectors to dense 100-dimensional representations,
        reducing parameter count by $\approx 750\times$.

  \item \textbf{Forward Pass Mechanics.}  The RNN processes the embedded
        sequence token by token.  After the final token, the hidden
        state $h_T \in \mathbb{R}^{64}$ encodes the entire sequence
        and is fed to the linear layer for binary classification.

  \item \textbf{Training Loop.}  \texttt{BCEWithLogitsLoss} +
        Adam optimiser; 10 epochs on 5\,000 samples.  Loss decreased
        monotonically, confirming correct implementation.

  \item \textbf{Vanishing Gradient and LSTM Motivation.}  The unrolled
        RNN at depth $T$ suffers from vanishing gradients because
        gradient products shrink exponentially.  LSTM addresses this
        with a gated cell state; GRU uses a simplified two-gate design.
\end{enumerate}

\subsection{Key Equations}
\label{subsec:key_equations}

For quick reference, the three core equations from this session:

\begin{align}
  \text{Hidden state update:} &\quad
    h_t = \tanh(W h_{t-1} + U x_t + b)
    \label{eq:hidden_summary} \\
  \text{Output at } t=3: &\quad
    O_3 = f(W \cdot O_2 + U \cdot X_3)
    \label{eq:o3_summary} \\
  \text{Binary cross-entropy loss:} &\quad
    \mathcal{L} = -\bigl[y \log \sigma(\hat{y})
                       + (1-y)\log(1-\sigma(\hat{y}))\bigr]
    \label{eq:bce_summary}
\end{align}

\subsection{Looking Ahead}
\label{subsec:ahead}

The next session will introduce \textbf{Long Short-Term Memory (LSTM)}
networks.  Key topics to anticipate:

\begin{itemize}[noitemsep]
  \item The cell state $c_t$ as a ``conveyor belt'' for long-range information.
  \item The forget gate, input gate, and output gate.
  \item How LSTM mitigates the vanishing gradient problem.
  \item PyTorch implementation using \texttt{nn.LSTM}.
\end{itemize}

Students were encouraged to practise by:
\begin{enumerate}[noitemsep,label=(\roman*)]
  \item Printing the output of every layer in the notebook to verify
        shapes at each stage.
  \item Extending the training to the full 50\,000-sample dataset and
        comparing accuracy.
  \item Replacing the embedding with pre-trained GloVe vectors and
        observing any accuracy improvement.
\end{enumerate}

\begin{keyinsight}[title=The Big Picture]
  The RNN family -- from the vanilla network in this session through
  LSTM and GRU -- all share the fundamental design principle of
  \cref{eq:rnn_update}: a hidden state that is updated at each time
  step by combining the current input with the previous hidden state.
  Every subsequent innovation (LSTM gates, attention mechanisms,
  Transformers) can be understood as a solution to specific limitations
  of this basic recurrence.
\end{keyinsight}

%%%%%%%%%%%%%%%%%%%%%%%%%%%%%%%%%%%%%%%%%%%%%%%%%%%%%%%%%%%%%%%%%%%%%%%%%%%%%
%% References / Resources
%%%%%%%%%%%%%%%%%%%%%%%%%%%%%%%%%%%%%%%%%%%%%%%%%%%%%%%%%%%%%%%%%%%%%%%%%%%%%
\newpage
\section*{Session Resources}
\addcontentsline{toc}{section}{Session Resources}

\begin{itemize}
  \item \textbf{Video lecture:}
        \url{https://vimeo.com/1156130746}
  \item \textbf{Slide deck (PDF):}
        \url{https://learning.iiitdwd.ac.in/mod/pdfprotect/view.php?id=1556}
  \item \textbf{Notebook:} \texttt{RNNClassification.ipynb}
        (shared via course portal; PyTorch implementation of the
        IMDB sentiment classifier)
  \item \textbf{Dataset -- IMDB:}
        \texttt{datasets.load\_dataset("imdb")} via HuggingFace, or
        download from Kaggle
  \item \textbf{Dataset -- Shakespeare:}
        \texttt{kylehobbsdev/shakespeare-dataset-of-500-shakespeare-plays}
        on Kaggle Hub
\end{itemize}

\end{document}
